\chapter{L4 — Instances of the computability structure}

% P-R reference: Chapter 2, examples 2.1, 2.4-2.5, literature/papers/PourEl-Richards-chapt2.md.
% Layer status: C[a,b] is stub (A1/A2/A3 theorem-body sorries grouped at
% ComputableAnalysis/L4/Instances/CMap.lean:261). L^p[a,b] and separable Hilbert
% are pending.

\begin{theorem}
  \label{thm:l4_cmap_instance}
  \lean{ComputableAnalysis.L4.instComputabilityStructureCMap}
  \uses{def:l3_computabilityStructure}
  $C([a,b], \mathbb{R})$ with sup-norm carries a \texttt{ComputabilityStructure}~$\mathbb{R}$ instance via the polynomial-approximation predicate. Stated in Lean; three axiom fields (\texttt{axiom\_linearity}, \texttt{axiom\_limits}, \texttt{axiom\_norms}) currently \texttt{sorry}. P-R Ch.\,2:141--153.
  See \texttt{ComputableAnalysis/L4/Instances/CMap.lean:261}.
\end{theorem}

\begin{theorem}
  \label{thm:l4_lp_instance}
  \notready
  $L^p[a,b]$ (for $1 \le p < \infty$) carries a \texttt{ComputabilityStructure}~$\mathbb{R}$ instance via $L^p$-norm-computable rational-step-function approximations. P-R Ch.\,2:146.
\end{theorem}

\begin{theorem}
  \label{thm:l4_hilbert_instance}
  \notready
  Every separable Hilbert space $E$ over $\mathbb{R}$ carries a \texttt{ComputabilityStructure}~$\mathbb{R}$ instance via a computable orthonormal basis. P-R Ch.\,2:147--148.
\end{theorem}
